\documentclass[11pt]{article}
\usepackage[a4paper,margin=24mm,headheight=15pt]{geometry}
\usepackage[T1]{fontenc}
\usepackage{lmodern,microtype,amsmath,amssymb,amsthm,booktabs,longtable,array,enumitem,fancyhdr,xurl}
\usepackage[hidelinks]{hyperref}
\newtheorem{theorem}{Theorem}[section]
\newtheorem{lemma}[theorem]{Lemma}
\newcommand{\Con}{\operatorname{Con}}
\newcommand{\file}[1]{{\small\path{#1}}}
\setlength{\parindent}{0pt}\setlength{\parskip}{5pt}
\setlist{nosep,leftmargin=1.6em,topsep=3pt}
\pagestyle{fancy}\fancyhf{}\fancyhead[L]{\small K215: a different encoding and counter seed}\fancyhead[R]{\small 21 October 2026}\fancyfoot[R]{\thepage}
\hypersetup{pdftitle={Two alternative 215-state ZFC consistency machines},pdfauthor={Research contribution for the ProveIt project}}
\begin{document}
\begin{center}
{\LARGE\bfseries Two alternative 215-state\par ZFC consistency machines}\par\vspace{8pt}
{\large Re-encoding the source and admitting a nonzero initial counter}\par\vspace{5pt}
{Research contribution for Vladimir Reshetnikov's ProveIt project}\par
{21 October 2026 UTC}
\end{center}
\begin{abstract}
We give two binary, single-tape, blank-start machines with 215 working states whose nonhalting is equivalent to the consistency of ZFC. The primary construction is a source re-encoding. This is an alternative construction, not a 214-state result. Its transition graph is not a relabeling of the three compared 215-state baselines. Relative to the primary baseline, it changes 17 transition rows and 4,519 source instruction addresses, while retaining the same 4,751 primitive operations and their complete guarded ordinal control graph. The number of addresses containing six consecutive zero bits decreases from 37 to 33; the source DAG remains at 202 nodes.

The important additional proof obligation is initialization. The new machine reaches its genuine source root after 716,307 transitions with next-candidate register 4 and additive counter 1, rather than the old zero-counter seed. We prove that an arbitrary natural initial counter is harmless: command soundness already permits it, and a cofinal family of refutation programs begins by resetting it. Fresh unbounded runtime verification establishes all 6,176 guarded primitive cases for this literal table. The resulting proof is ordinary mathematics formalizable in ZFC, not a newly kernel-checked Lean or Rocq theorem. The calculus itself is inherited unchanged; the primary contribution is a materially different physical encoding and a more flexible initialization theorem. A second table gives a distinct physical allocation of the recently published sixth-axiom sampler source, together with its adjusted strong compiler and separate runtime and blank-start verification.
\end{abstract}
\setcounter{tocdepth}{1}\tableofcontents
\section{Result, distinction, and scope}
The machine uses one bi-infinite binary tape and moves left or right at every step. It has 215 working states and 430 transition rows. Including the halt label gives 216 labels, not 216 working states. All working states are reachable in the directed transition graph; this is not a claim that every state occurs on the blank computation.

The desired conclusion is
\[
\mathrm{ZFC}\vdash
\bigl(\operatorname{Halts}(M,\mathrm{blank})
\leftrightarrow\exists p\,\operatorname{Proof}_{\mathrm{ZFC}}(p,\bot)\bigr).
\]
Equivalently, ZFC proves that blank nonhalting of $M$ is equivalent to $\Con(\mathrm{ZFC})$. No assertion of actual nonhalting, consistency of ZFC, minimality, or priority over every unpublished construction is made.

\subsection{What is materially different?}
The primary reference is the earlier composed-command/partial-selector K215 construction. A second reference permutes the DROP and NIP tags. A recent public variant exchanges two helper roles, at native labels 11 and 134. The new table is not isomorphic to any of these as a transition graph preserving the input symbol, written symbol, and move direction, even when the initial state is not distinguished. Joint color refinement distinguishes the respective graph multisets in 3, 3, and 2 rounds. This is a finite certificate of nonisomorphism, not a general graph-isomorphism heuristic: unequal refinement multisets preclude an isomorphism.

Against the primary reference, 17 retained transition rows change and 4,519 of the 4,751 primitive addresses change. Against the other two references, the literal native-label differences comprise 33 and 22 rows respectively. Native-label differences alone would not prove material novelty; the graph certificate and different blank computation provide the stronger evidence.

The old primary reaches a clean $(N,D)=(4,0)$ source root after 95,714 steps. The new table reaches $(N,D)=(4,1)$ after 716,307 steps. These are different physical computations and different proof boundaries. The old initializer certificate cannot simply be copied to the new table.

The logical calculus and abstract register program are unchanged. In particular, this is not a new axiomatization of set theory, a new command calculus, or a further reduction in state count. The change is a different allocation of an equivalent source encoding, plus a proof that a less restrictive initialized counter is sufficient.

\subsection{Pinned artifacts}
The native transition table has SHA-256
\begin{quote}\small\path{f5acfbb2e3d0e2b61350454a098a217d3457710311e9dbf02135098fdbe8d36d}\end{quote}
The normalized discovery table has SHA-256
\begin{quote}\small\path{6aab0397336a198c8fd765803864a69134508175e637d563d6bea8aabbe7b416}\end{quote}
The exact indented source-site file has SHA-256
\begin{quote}\small\path{704a0449c4164c39464f714e08c435d656f96c024e336ff27d091166772816fe}\end{quote}
These hashes identify literal files, including their serialization. The native, normalized, and dense canonical tables are explicitly bijective relabelings of this one construction. The machine definition is the table, not its name or these state-label conventions.

\section{Equivalent source encoding and physical verification}
\label{sec:source-correspondence}
The source has 4,751 primitive occurrences and 202 distinct subtrees. Reassociation changes source addresses and the corresponding physical routing, while preserving the literal left-to-right operation/operand word and both guarded successor choices at every conditional primitive. The first 37 primitive addresses remain fixed. The maximum address depth increases from 47 to 49. The number of addresses containing a run of six zero bits decreases from 37 to 33. This last quantity counts addresses, not states or individual zero runs, and is not claimed to be a complexity lower bound.

\begin{lemma}[All-natural source correspondence]
Relate the old and new source configurations when their primitive ordinals and all nine natural-valued registers agree. This relation is a strong step-for-step bisimulation, preserving halt.
\end{lemma}
\begin{proof}
At each ordinal the primitive operation and operand are identical. Therefore its arithmetic update and guard outcome agree on every natural store. The complete guarded ordinal control graphs agree, so the next ordinal or halting result agrees as well. Induction on execution length gives the claim. This argument has no bounded-register assumption and no formula-decoding premise.
\end{proof}

Source correspondence does not by itself prove that the reallocated binary table implements that source. A separate symbolic runtime check treats every primitive occurrence, splitting decrement cases according to zero or positive input. The fresh check proves all 6,176 cases and uses 1,753,058 symbolic rules. The minimum constant term in the certified affine transition duration is 6; all duration coefficients are nonnegative. Thus every source primitive corresponds to a positive finite binary execution, ruling out a zero-time simulation loophole.

The representation contracts allow arbitrary natural primary registers, an arbitrary natural padding gap, and a nonempty finite spare list, together with a completely zero exterior. Scan and list lemmas are inductive, not bounded test traces. Address decoding, both guarded exits, the represented output values, and the tape outside the represented frame remain separate proof obligations. The supplied proof machinery rechecks these rather than inferring them from the source DAG size.

The source/calculus bridge is inherited from the earlier composed K215 and its local-input K216 parent. The following sections reproduce the numerical calculus and explicit command compiler to make clear exactly which mathematics is retained. The final seed theorem extends their stated initialization boundary; it does not silently identify the new counter value with zero.

\section{The numerical calculus and its consistency calibration}
\label{sec:native-calibration}

The present construction uses the dominated-parameter calculus supplied in
\file{K218_Dominated_Rho8_2026-10-20.zip}.  This is a recent input to the
repository, distinct from its primary, unchanged-calculus K218 table.
Its logical theorem is inherited through the 217-state parent. The present construction retains that numerical calculus and changes the physical source encoding and blank initialization.  This distinction is essential for identifying
which argument supplies each part of the result.

\subsection{Raw numerical operations}

For natural numbers put
\[
 P(a,b)=\frac{(a+b)(a+b+1)}2+b,
 \qquad I(a,b)=4P(a,b),\qquad a\to b:=I(a,b).
\]
The remaining constructors are
\[
 R(a,b)=4P(a,b)+1,\quad Q_a b=4P(a,b)+2,
 \quad E(a,b)=4P(a,b)+3.
\]
Raw truth is $T=0$ and the refutation target is $F=1=R(0,0)$.
Write $NX=I(X,F)$ and $\exists_a X=N(Q_a(NX))$.  In particular,
$I(0,0)=0$: zero is simultaneously the designated initial theorem and
an implication code.  All quantified schema parameters below range over
\emph{every natural number}.  No test that a number decodes as a conventional
formula is introduced.

Let $\mathcal C$ be the accepted native calculus, called \texttt{C222} in
the inherited interfaces and used by K219.  It has truth, numerical modus
ponens, and theorem-only generalization:
\[
 \frac{\mathcal C\vdash X\quad \mathcal C\vdash I(X,Y)}
      {\mathcal C\vdash Y},
 \qquad
 \frac{\mathcal C\vdash X}{\mathcal C\vdash Q_aX}.
\]
Its eight logical rows comprise the following seven retained families
and the unrestricted family $L$ displayed below:
\begin{align*}
 \rho_1(a,b)&=\exists_a E(a,b),\\
 \rho_2(a,b,c)&=Q_a(b\to c)\to(Q_ab\to Q_ac),\\
 \rho_3(a,b,c)&=E(a,b)\to(R(c,a)\to R(c,b)),\\
 \rho_4(a,b,c)&=Q_aQ_bc\to Q_bQ_ac,\\
 \rho_5(a,b,c)&=E(c,a)\to(R(c,b)\to R(a,b)),\\
 \rho_6(a,b)&=(\exists_a Q_ab)\to b,\\
 G(a,b,c)&=E(c,a)\to\bigl(E(c,b)\to Q_{R(a,b)}E(a,b)\bigr),\\
 L(A,B,C,U)&=((C\to U)\to A)\to((A\to C)\to(B\to C)).
\end{align*}

The accepted set-theoretic information remains in the literal gated seed
bundle.  To specify this inherited object, define
\begin{align*}
 B(b,k,m)&=Q_b\bigl(m\to NQ_b(m\to NR(b,k))\bigr),\\
 S(a,b,k)&=Q_k\bigl(R(k,a)\to R(k,b)\bigr),\\
 C(a,b,k,m)&=Q_b\bigl(R(b,k)\to
           NQ_a(m\to NQ_a(S(a,b,k)\to B(b,k,m)))\bigr),\\
 J(a,b,c)&=Q_c\bigl(C(a,b,c,c)\to NB(b,c,c)\bigr),\\
 V(a,b,c)&=R(a,b)\to Q_cR(a,b),\\
 \operatorname{Bundle}(a,b,c)&=N\bigl(V(a,b,c)\to J(a,b,c)\bigr).
\end{align*}
A bundle is an axiom when $c\ne a$ and $c\ne b$.  Both occurrences of
$c$ as matrix code and private label are retained literally.  We neither
remove these gates nor replace the bundle by an informal set-theoretic
analogue.  The source returns truth instead when a gate fails.

The established native calibration supplies finite-refutation
transformations between $\mathcal C\vdash F$ and ZFC inconsistency.
That inherited result is a premise of the present construction.  It is
supported by the repository's native-calculus developments and their
connection to the compact set theory underlying K230; it is not proved
anew by a transition-table simulation.

For provenance, the dominated package pins its K219 inputs to ProveIt
revision \path{f8ec6158f3bdc885139ee395ae34fe44c3251b9f} and includes
its exact native interfaces and finite logic certificate.  Earlier proof-assistant endpoints concern their own exact tables. The
native-calculus bridges, dominated theorem and present source/runtime
arguments must all be included to reach the new table. No new exact-215
proof-assistant endpoint or internal PA proof object is asserted here.

\subsection{Why a dominated fourth parameter is sufficient}

Define $\mathcal D$ by retaining all of the preceding clauses, except
that its $L(A,B,C,U)$ axiom requires
\[
                         U\le A.                 \tag{D}
\]
This is an ordinary inequality between natural codes, not a syntactic
freshness test and not an inequality of semantic truth values.

\begin{theorem}[Inherited dominated-parameter theorem]
For every natural code $n$,
\[
                  \mathcal D\vdash n
             \quad\Longleftrightarrow\quad
                  \mathcal C\vdash n.
\]
The proof is uniform in arbitrary native codes, including zero and
all degenerate constructor coincidences.
\end{theorem}

The easy direction follows because every dominated instance is already
an unrestricted instance.  For the converse, the inherited finite
certificate derives the entire unrestricted $L$ schema from truth,
modus ponens, and dominated instances alone.  Here are the arithmetic
and proof-theoretic facts that make that certificate meaningful.

First, $P(a,b)\ge a,b$ for all naturals.  If $a+b=s\ge1$, then
$s(s+1)/2\ge s$, and the case $s=0$ is immediate.  Consequently
$I(a,b)\ge a,b$.  If a schematic term $U$ is literally a subterm of an
implication term $A$, then $U\le A$ under every numerical assignment.
The same conclusion holds whenever $U=T$.  This remains true when some
implications evaluate to zero.  The finite checker uses precisely these
sufficient conditions for all dominated leaves.

Second, every certificate node denotes a universally quantified
\emph{numerical} statement.  Replacing one of its metavariables by an
implication term is therefore universal instantiation.  It does not
assume a general substitution rule inside the object calculus, and it
does not rename the numerical labels of a previously decoded atom.
At an MP node, equality of the instantiated antecedent and minor is
checked exactly.  The sole native normalization is $I(T,T)=T$.

The unrestricted-row certificate has the following exact inventory:
\begin{center}
\begin{tabular}{lr}
\hline
Item & Count\\
\hline
Proof nodes & 3,728\\
Shared term nodes & 5,160\\
Dominated-$L$ leaves & 1,260\\
Truth leaves & 1\\
MP steps with explicit substitutions & 55\\
Direct universal-instantiation steps & 410\\
Direct MP steps & 2,002\\
Maximum checked subterm depth & 4\\
\hline
\end{tabular}
\end{center}
These are finite proof objects, not samples of numerical assignments.
The implication-subterm inequality proves each axiom side condition for
all assignments; induction through the checked proof proves the
schematic result for all assignments.

The certificate first obtains identity, weakening $K$, Peirce's law,
composition, and distribution $S$:
\begin{align*}
 K&:p\to(q\to p),\\
 \mathrm P&:((p\to q)\to p)\to p,\\
 \mathrm B&:(p\to q)\to((q\to r)\to(p\to r)),\\
 S&:(p\to(q\to r))\to((p\to q)\to(p\to r)).
\end{align*}
A key intermediate theorem is
\[
 (p\to q)\to\bigl(((p\to r)\to p)\to q\bigr).
\]
The successful inherited proof of that theorem adds the dominated
context $A=((z\to U)\to v)\to w$; the fourth argument lies three
implication edges inside the first.  The rest of the fixed derivation
reaches unrestricted $L$.  The complete certificate and its checkers
are included with the inherited dependency, so the large intermediate
formulas need not be expanded into the body of this article.

For clarity, the final passage from $K,S,\mathrm P$ to $L$ has a short
ordinary explanation.  Temporarily assume
$H=(C\to U)\to A$ and $J=A\to C$.  Under the additional assumption
$C\to U$, two MPs give $A$ and then $C$.  The deduction transformation
justified by $K,S$ gives $(C\to U)\to C$.  Peirce's law at $C,U$
gives $C$.  Weakening gives $B\to C$, and discharging $J,H$ yields
exactly $L(A,B,C,U)$.  No generalization is used under these temporary
assumptions.  The supplied certificate expands this deduction argument
into explicit schema instances and MP steps.

Replacing each unrestricted-$L$ leaf of a finite $\mathcal C$ proof by
this fixed derived proof gives a finite $\mathcal D$ proof of the same
code.  All other axioms and both rules are retained literally, proving
the theorem.  Combining it with the inherited native calibration gives
\[
       \mathcal D\vdash1
       \quad\Longleftrightarrow\quad
       \text{there is a finite ZFC refutation}.       \tag{C1}
\]


% Include-only supplement. Requires amsmath and amssymb.
% The finite source checker is audit_command_permutation.py.
\begingroup
\newcommand{\Kp}{\mathsf P}
\newcommand{\Ki}{\mathsf I}
\newcommand{\Kq}{\mathsf Q}
\newcommand{\Kd}{\mathcal D}
\newcommand{\Kax}{\mathsf{AX}}
\newcommand{\Kmp}{\mathsf{MP}}
\newcommand{\Kdrop}{\mathsf{DROP}}
\newcommand{\Knip}{\mathsf{NIP}}
\newcommand{\Kenc}{\operatorname{enc}}
\newcommand{\Kblk}{\mathcal B}
\newcommand{\Kdiscard}{\mathsf{Discard}}

\section{Inherited all-natural calibration of the local-input command source}
\label{sec:k216-source-calibration}

This section reproduces, with attribution, the explicit ordinary calibration argument supplied with \file{K216_Permuted_Commands_2026-10-21.zip}, with attribution to that construction. Its source hashes and selector paths identify the parent reference trees, not the new reparenthesized source. The exact all-natural source correspondence in Section~\ref{sec:source-correspondence} transports this theorem to the present source without changing any arithmetic command or candidate encoding.

The inherited source-level transformation is a
permutation of the three small command selectors. The new construction keeps
these selectors; its additional initialization theorem is proved below. The arithmetic
commands and the dominated axiom family are unchanged. This section
proves that the permutation preserves the full consistency calibration,
including raw candidate numbers, the zero alias, arbitrary formula
codes, and the actual source seed with next-candidate register $4$.
It does not infer a binary runtime theorem from a source subtree count;
the binary runtime and blank-tape certificates are separate obligations.

\subsection{Exact finite source identity}

Write
\[
 \Kp(a,b)=\frac{(a+b)(a+b+1)}2+b,\qquad
 \Ki(a,b)=4\Kp(a,b),\qquad
 \Kq_i a=4\Kp(i,a)+2.
\]
The native truth and falsity codes are respectively $0$ and $1$.
In particular, $\Ki(0,0)=0$. Let $\Kd$ be the dominated numerical
calculus of the preceding dominated-rho8 construction: truth, numerical
modus ponens, theorem-only generalization, the seven retained logical
rows, the gated theory bundle, and
\[
 L(A,B,C,U)=
 \Ki\bigl(\Ki(\Ki(C,U),A),
            \Ki(\Ki(A,C),\Ki(B,C))\bigr),\qquad U\leq A.
\]
The inherited finite propositional certificate proves equality of the
$\Kd$ and K219 theorem sets for \emph{every natural code}. Together with
the inherited numerical-calculus-to-ZFC calibration, it gives
\begin{equation}
 \Kd\vdash1\quad\Longleftrightarrow\quad
 \exists p\;\operatorname{Proof}_{\mathrm{ZFC}}(p,\bot).
 \label{eq:k216-inherited-calibration}
\end{equation}
No well-formedness or successful formula-decoding test is inserted into
this equivalence.

The source checker first reconstructs the 205-subtree noncommon-input
reference directly from the pinned K219 source. Its independently
implemented structured-control and address-counter interpreters agree
on every primitive and every branch outcome. The resulting guarded
ordinal control graph is exactly that of the frozen 202-subtree source
with SHA-256
\begin{center}\ttfamily\footnotesize
abb7229877759170b61dfcfe21ca0889b3df2f91f4215383817a125eb0f8996c.
\end{center}
Thus the frozen source has the arithmetic command semantics proved in
the dominated-rho8 report, with counter/cache field order YX.

In this frozen source, the MP, DROP, and NIP bodies occur at paths
$111$, $110101$, and $110100101$, respectively. They contain
$106$, $26$, and $26$ primitive occurrences. Each has zero escaping
break and continuation demand. NIP contains the halting primitive;
all its nonhalting exits are normal returns. The branch-skipping
instructions outside the three bodies remain fixed when the bodies
are permuted. The checker traces their ordinary continuations: the
last body returns directly to source ordinal $0$, the middle body
returns through the unconditional break at $4724$, and the innermost
body returns through the breaks at $4697$ and $4724$. Each route ends
at ordinal $0$ without another arithmetic command.

For the source denoted $210$, put NIP, DROP, and MP in these three
selector positions. For $120$, put DROP, NIP, and MP there instead.
The exact selector partitions and wire formats are
\[
\begin{array}{c|c|c|c|c}
 &\text{selector }0&\text{selector }1&\text{selector }2
 &\text{selector }c+3\\\hline
210&\Knip&\Kdrop&\Kmp&\Kax(c;x,y)\\
120&\Kdrop&\Knip&\Kmp&\Kax(c;x,y)
\end{array}
\]
and
\begin{align*}
 f_{210}(\Knip)&=[0],& f_{210}(\Kdrop)&=[1],&
 f_{210}(\Kmp(x))&=[2,x],\\
 f_{120}(\Knip)&=[1],& f_{120}(\Kdrop)&=[0],&
 f_{120}(\Kmp(x))&=[2,x],\\
 f_\pi(\Kax(c;x,y))&=[c+3,y,x]&& (\pi=210,120).
\end{align*}
The semicolon in $\Kax(c;x,y)$ separates the selector parameter from
the counter increment $x$ and cache increment $y$; it does not denote
wire order.

The independent checker proves these partitions directly from the three
selector decrement tests, at primitive ordinals $23,24,25$. In the
$210$ source, their zero branches lead to NIP at ordinal $4725$, DROP
at $4698$, and MP at $4591$. Three positive outcomes enter AX at ordinal
$26$, with residual selector $c$. This last case is the symbolic case
of every input selector at least $3$, not a finite sampling of large
selectors.

Both sources have 202 distinct subtrees, 4,751 primitive occurrences,
6,176 guarded primitive cases, maximum address depth 44, and maximum
zero run 6. All 4,593 primitive occurrences exterior to the replaced
bodies, including the entire AX subtree, remain literal. The checker
also verifies that each candidate tree is \emph{exactly} the indicated
permutation of the frozen tree. The candidate source hashes are
\begin{center}\footnotesize\ttfamily
210: 8b5dc46f10d927b2db5e4bcb522d482779ec61c72c82717d7bf12fca3550a3bc\\
120: ee1f7c7cb216a80315449371b040ee094728a655439954225ed4be4635296680.
\end{center}

\subsection{Command equations and soundness for raw codes}

At a command boundary write the retained tuple as $(t,w,d,h)$: top,
encoded stack, additive counter, and cached formula. The notation
$[x_1,\ldots,x_k;w]$ means right-nested pairing, and
$\operatorname{unpair}(0)=(0,0)$. The exact command equations, independent
of which small tag names the command, are
\begin{align*}
 \Kmp(x)(t,w,d,h)&=
 \begin{cases}
 (d+x,\Kp(0,w),d+x,0),&\Ki(t,d+x)=h,\\
 (0,w,d+x,0),&\Ki(t,d+x)\ne h;
 \end{cases}\\
 \Kdrop(t,w,d,h)&=(h,\Kp(t,w),d,h),\\
 \Knip(t,w,d,h)&=
 \begin{cases}
 \mathsf{halt},&h=1,\\
 (t,w',d,h'),&(h',w')=\operatorname{unpair}(w),\ h\ne1.
 \end{cases}
\end{align*}
For $\Kax(c;x,y)$ put
\[
 U=d,\qquad A=d+x,\qquad B=h+y,\qquad C=c,
\]
and let $G(A,B,c)$ be zero when $c=A$ or $c=B$, and the retained theory
bundle otherwise. In the inherited stack order its output is
\begin{align*}
 t'&=G(A,B,c),& d'&=0,&h'&=0,\\
 w'&=[R_1(A,B,c),\ldots,R_7(A,B,c),L(A,B,c,U),\Kq_Uh,t;w].
\end{align*}
Here the $h$ in $\Kq_Uh$ is the incoming cache, before adding $y$.
The bound $U\leq A$ holds for every raw payload because $A=d+x$.

A body with no escaping break or continuation has the same effect at
any of the three command positions. The branch wrappers route every
normal return back to the command loop. Consequently the equations
above hold for the permuted sources on every natural input tuple,
including payloads arising from incomplete or otherwise unintended
candidate encodings.

\paragraph{Soundness invariant.}
If $t$, $h$, and every successive head of $w$ are $\Kd$-theorems,
every nonhalting command preserves that property. A successful MP is
exact numerical modus ponens; a failed MP inserts truth $0$. DROP and
NIP only move theorem codes. AX inserts retained axioms, a valid
dominated-row instance, a generalization of a theorem, and either a
permitted bundle or truth. The only halt is the NIP test $h=1$.
Therefore halting from any such safe seed implies $\Kd\vdash1$.
This argument quantifies over all natural codes and does not discard
``junk'' candidates.

\subsection{An explicit compiler retained under the tag change}

For completeness, a strong block $\Kblk_X$ must, from every $(t,w,0,0)$,
either halt or return $(X,w,0,0)$. The command lists below use command
\emph{names}; their concrete wire encodings use $f_{210}$ or $f_{120}$.
Concatenation denotes sequential execution, and a superscript on a
command denotes repetition.

Define
\[
 \Kdiscard=[\Kmp(1),\Kax(1;0,0)]\Knip^{10},\qquad
 J(u)=\begin{cases}[],&u=0,\\{}[\Kmp(u)],&u>0.\end{cases}
\]
The first MP fails because $\Ki(t,1)>0$, whereas the incoming cache is
zero. The following AX has $U=A=c=1$ and a zero gate. Its nine newly
generated codes are all different from $1$, and its saved old top is
zero. Hence ten NIPs return exactly $(0,w,0,0)$.
The block $J(u)$ maps $(0,w,0,0)$ to $(0,w,u,0)$: if $u>0$,
$\Ki(0,u)>0$ forces failure. Omitting MP when $u=0$ avoids the extra
stack zero that a successful $\Kmp(0)$ would insert.

Truth uses $\Kdiscard$. A dominated-row leaf with $U\leq A$ uses
\[
 \Kdiscard\;J(U)\;[\Kax(c;A-U,B)]\;\Knip^8\;[\Kdrop]\;\Knip^3.
\]
For retained row $R_j(A,B,c)$, where $1\leq j\leq7$, use
\[
 \Kdiscard\;[\Kax(c;A,B)]\;\Knip^j\;[\Kdrop]\;\Knip^{11-j}.
\]
For a permitted bundle use
$\Kdiscard[\Kax(c;A,B)]\Knip^{10}$.
These suffixes select the indicated top or stack entry and restore the
underlying stack and zero cache. After expanding the negation and
existential abbreviations, all eight row codes are implications,
and $\Kq_Uh\equiv2\pmod4$, so none of these generated stack codes is
falsity $1$.

Given a block for $X$, a block for $\Kq_iX$ is
\[
 \Kblk_X\;[\Kdrop]\;J(i)\;[\Knip,\Kax(i;0,0)]\;
 \Knip^9\;[\Kdrop,\Knip,\Knip].
\]
After DROP and $J(i)$, the current top and cache are zero, the counter
is $i$, and $X$ is the first stack head. NIP places $X$ in the cache.
AX then has $U=A=c=i$, so its top gate is zero and its ninth generated
stack entry is exactly $\Kq_iX$. The selection suffix returns that code
with the original stack and zero counter and cache.

Given blocks for $X$ and $\Ki(X,Y)$, use
\begin{multline*}
 \Kblk_X\;[\Kdrop]\;\Kblk_{\Ki(X,Y)}\;
 [\Knip,\Kdrop,\Knip,\Kmp(Y),\Kax(Y;0,0)]\\
 {}\Knip^{10}\;[\Kdrop,\Knip,\Knip].
\end{multline*}
Unless an earlier NIP has halted, the two premises reach the top and
cache in the required order, and MP succeeds. The reset AX has
$U=A=c=Y$ and a zero gate. The final selection returns $Y$ with the
original stack and zero counter and cache, or halts if a tested cache
is already $1$.

Induction on a finite $\Kd$ derivation proves the strong block contract
for every derived natural code. The construction is primitive recursive
on the derivation. The new selectors alter only the wire names of these
commands, so every calculation above applies verbatim to both permuted
sources.

\subsection{Exact packing and cofinality from the actual cutoff}

For a finite field list $fs$ and tail $r$, write $E(fs,r)$ for its
right-folded Cantor pairing. Command tails are encoded by
\[
 \Kenc_\pi([])=0,\qquad
 \Kenc_\pi(k::L)=1+E(f_\pi(k),\Kenc_\pi(L)).
\]
A newly enumerated candidate is loaded directly from its number. A
subsequent command is decoded from one less than the current positive
proof-list value. Thus the number for a whole first candidate
$k::L$ is $E(f_\pi(k),\Kenc_\pi(L))$, without an outer successor.
This distinction is essential when NIP has tag zero.

Every field reader returns a tail no greater than its input. Therefore
one command starting from a positive proof-list value $v$ returns a
proof-list value at most $v-1$, unless it halts. All arithmetic macros
terminate on all natural inputs. The full macro contracts also restore
selector $A=0$ and scratches $X=Y=0$ at every nonhalting command return;
they preserve $N$ outside candidate fetch and leave $L$ equal to the
unconsumed unpairing remainder. At fresh fetch the source loads $N$
as the new candidate and increments $N$ exactly once. Every nonhalting
candidate therefore drains after finitely many commands; the search
then proceeds to the next candidate. In particular, zero selectors or
zero field tails do not cause an omitted case or a nonhalting inner
computation.

Suppose $\Kd\vdash1$, let $\Kblk_1$ be its compiled block, and put
\[
 S=[\Kdrop,\Knip,\Knip],\qquad
 R=\Kenc_\pi(\Kblk_1S).
\]
For any natural padding value $Y$, prepend
$\operatorname{Start}_Y=\Kax(0;1,Y)$. It sets the counter and cache to
zero from every incoming retained tuple. Thus
$[\operatorname{Start}_Y]\Kblk_1S$ halts from every incoming tuple:
the block either halts, or returns top $1$ and zero cache, and $S$
then puts $1$ in the cache and tests it.
Its exact first-candidate number is
\begin{equation}
 N_Y=\Kp\bigl(3,\Kp(Y,\Kp(1,R))\bigr)\geq Y.
 \label{eq:k216-cofinal}
\end{equation}
The inequality follows from the coordinate inequalities for Cantor
pairing. This gives a cofinal family of successful candidates for both
permutations. Choose $Y\geq4$. The successful candidate $N_Y$ is at or
beyond the actual initial next-candidate value $4$. Finite draining of
earlier candidates guarantees that the search reaches it unless a halt
has already occurred. More generally the same proof works beyond every
natural cutoff, and after every finite initial partial candidate.

\paragraph{Source calibration theorem.}
For either source permutation, started at its genuine source root with
register tuple
\[
 (N,T,W,A,D,H,X,Y,L)=(4,0,0,0,0,0,0,0,0),
\]
ZFC proves
\[
 \operatorname{Halt}_{\rm source}
 \quad\Longleftrightarrow\quad\Kd\vdash1
 \quad\Longleftrightarrow\quad
 \exists p\;\operatorname{Proof}_{\mathrm{ZFC}}(p,\bot).
\]
\emph{Proof.} The zero top, zero cache, and all heads of the zero stack
are truth theorems. The soundness invariant proves the forward
implication. The explicit compiler, cofinal family
\eqref{eq:k216-cofinal}, and draining argument prove the reverse
implication. Equation~\eqref{eq:k216-inherited-calibration} gives the
last equivalence. The argument uses finite source identities,
primitive-recursive proof and word transformations, and induction on
natural numbers and finite derivations. It is formalizable in ZFC
without assuming $\operatorname{Con}(\mathrm{ZFC})$.

Combining this theorem with a separate exact binary runtime simulation
and a blank-tape trace to the displayed represented source root yields
\[
 \mathrm{ZFC}\vdash
 \bigl(\text{the corresponding binary machine does not halt on blank tape}\bigr)
 \ \longleftrightarrow\ \operatorname{Con}(\mathrm{ZFC}).
\]
The trace's tape header, head position, gap and spare-word requirements
belong to that binary certificate; they are not consequences of the
source calibration alone.
\endgroup

\section{A generalized counter seed and the actual blank trace}
\label{sec:seed}
Write the nine logical registers in the order
$(N,T,W,A,D,H,X,Y,L)$. Here $N$ is the next candidate index, $T$ the current theorem, $W$ the encoded stack, $A$ the command selector, $D$ the additive counter, $H$ the theorem cache, $X,Y$ scratches, and $L$ the remaining proof-list code.

\begin{theorem}[Arbitrary-counter root calibration]
At a genuine source root, let $N=n$ and $D=d$ be arbitrary naturals, let $A=X=Y=L=0$, and suppose $T,H$ and all stack heads of $W$ are numerical theorems. Then source halting is equivalent to numerical inconsistency. In particular the theorem applies to $n=4,d=1,T=H=W=0$.
\end{theorem}
\begin{proof}
Soundness does not require $D=0$. In an AX command the old counter is $U=d$ and its new parameter is $A=d+x$, so the dominated-row requirement $U\le A$ holds for every natural payload $x$. The seven other axiom rows, theorem-only generalization, and the gated literal bundle retain their old soundness arguments. MP is numerical detachment when its equality test succeeds and produces truth otherwise. DROP and NIP only move theorem codes. The sole halting test is cached falsity $H=1$. Thus a halt implies inconsistency even when the initial counter is nonzero.

For the converse, take the strong compiled block $\mathcal B_1$ for a numerical derivation of falsity and append $S=[\mathsf{DROP},\mathsf{NIP},\mathsf{NIP}]$. The strong compiler works from every retained top and stack once $D=H=0$; it either halts or returns $(T,W,D,H)=(1,w,0,0)$ with the original stack $w$. The suffix then loads and tests falsity.

For each natural $y\ge1$, prefix this list by
$\mathsf{Start}_y=\mathsf{AX}(0;1,y)$. For arbitrary incoming $d,h$, this command has $U=d$, $A=d+1$, $B=h+y$, and $C=0$. Both gate inequalities hold because $A,B>0$. More importantly, irrespective of the incoming top, stack, or counter, the command sets $D=H=0$. Its new top and stack need not be prescribed: the strong block explicitly quantifies over arbitrary such data. Any earlier halt is already sufficient for the completeness direction.

Put $R=\operatorname{enc}_{210}(\mathcal B_1S)$. The exact first-candidate integer for the prefixed list is
\[
 n_y=P\bigl(3,P(y,P(1,R))\bigr)\ge y.
\]
There is no outer successor on a first candidate. Subsequent positive proof-list entries have the inherited successor convention. Choose $y\ge\max(1,n)$. Every preceding candidate either halts or drains in finite time: its first command performs finite arithmetic; each subsequent command strictly decreases its positive list remainder. The scheduler restores its selector and scratches and advances $N$. Consequently it reaches $n_y$ unless a halt has already occurred. That candidate resets the counter/cache, executes the compiled refutation, and halts. This proof does not assume a bound on the unknown counter carried into the successful candidate.
\end{proof}

\subsection{Literal blank initialization}
Exact simulation from the defined blank configuration reaches, after 716,307 transitions,
\[
 q=84,\quad h=-9,\quad g=152,\qquad
 (N,T,W,A,D,H,X,Y,L)=(4,0,0,0,1,0,0,0,0),
\]
in native labels and absolute tape coordinates. The spare list consists of 185 ones. The complete finite set of nonzero tape cells is supplied as a receipt and is checked, not merely the nine decoded numbers.

The leading one is exactly at $h-1=-10$. The next one is at $h+64+g=207$, so all 64 program-counter cells are zero. This is a genuine source root rather than an interior instruction with a conveniently decodable tail. The decoder checks the full unary-register frame, spare list, and zero exterior. No halt occurs in the verified prefix.

The zero current theorem and cache are truth; the zero stack has only truth under the inherited zero-pairing alias. The generalized theorem applies with $n=4,d=1$. Combining it with source bisimulation and the positive unbounded binary simulation proves the claimed blank-start consistency equivalence. All arguments concern finite tables, natural-number induction, primitive-recursive proof transformations, and finite derivations, and are formalizable in ZFC without assuming its consistency.

A separate reachability check contains the new counter-one seed in the old closed abstract invariant: its ternary class is 83 and its four-class index is 259. A fresh ternary closure reproduces the same bitmap. This is a useful cross-check, but it is not used as a substitute for either the arbitrary-counter theorem or the all-natural binary runtime proof.

\section{A second physical allocation with a different axiom sampler}
\label{sec:rho6-alternative}
A complementary alternative changes the \emph{source command effect}, rather than only its encoding. It emits $\rho_6(A,B)$ in place of $\rho_6(A,C)$, where
\[
 \rho_6(a,b)=(\exists_a Q_a b)\to b.
\]
The universal numerical calculus is still the same: only the way a command supplies an instance is changed. This sampler idea and its composed source are already recorded in public intake 199 at ProveIt commit
\file{7de71bd17b85db67529e8f10083c066ab0ed0c56}. Our contribution is a distinct physical allocation and its completed verification, not priority for the sampler idea.

Relative to the primary \emph{earlier} K215 table, not the counter-one table of this report, the new native table changes exactly two rows:
\[
\begin{array}{c|c|c}
\text{row}&\text{old action}&\text{new action}\\\hline
199/1&1R166&1R12\\
206/0&0R186&0R76
\end{array}
\]
No state is added or removed. Its native table SHA-256 is
\begin{quote}\small\path{a96732bd55c9412cbfd4bc3019200f364b3e3169f58e9a495dd98b072fb2d648}\end{quote}
Its canonical table SHA-256 is
\begin{quote}\small\path{bfa5ad02c53193656ed135588416bd05bdc72ffa2de92e3f6ed13556268df0d4}\end{quote} The four changed primitive operands are at zero-based ordinals $959,964,1137,1142$, changing register index 3 to 5. The source interval $[888,1171)$ is the 283-leaf sixth-axiom emitter. The tree retains 202 distinct subtrees and 4,751 leaves. Its exact source SHA-256 is
\begin{quote}\small\path{19da13e482ea0f3d57432aad852bdf2448093e58a6e3b3873ac744eae7440693}\end{quote}

\subsection{All-natural sampler theorem}
Conjugate the old emitter by swapping its parameter registers 3 and 5. The literal interval changes only the four listed CI/CD operands; all guarded successors remain fixed. The old emitter preserves its parameter registers and clears its scratches. Consequently the conjugate restores the same full input parameter store, constructs $\rho_6(A,B)$, and pushes the old current theorem on the same underlying stack. This is a stepwise all-natural register conjugacy, not an assumption that arbitrary derived formulas may be renamed or that a finite test set exhausts numerical inputs.

Every emitted code is therefore an instance of the same universally quantified native axiom. Other axiom emitters, the dominated fourth parameter, theorem-only generalization, and the literal gated bundle remain unchanged. The old theoremhood invariant still proves halting soundness for all raw codes.

Completeness needs an adjusted strong compiler for the sixth row. For arbitrary $a,b$, use
\[
 \mathsf{Discard};\mathsf{AX}(a;a,b);
 \mathsf{NIP}^{6};\mathsf{DROP};\mathsf{NIP}^{5}.
\]
After Discard, the top and counter/cache are zero and the underlying stack is the original $w$. This AX has $U=0,A=a,B=b,C=a$ and wire fields $[a+3,b,a]$. The bundle gate fails because $C=A$, so the new top is zero. After six pops, $\rho_6(a,b)$ is cached. DROP makes it current and pushes one extra zero. Five more pops remove that extra zero and the four remaining entries, ending with the saved old top zero in cache. The underlying stack is restored and $D=H=0$.

All intermediate axiom entries are implications, the generalized cache has tag 2, and the saved top is zero, so none is the refutation code 1. The other strong row blocks retain their selection positions and cleanup counts. Their unused sixth entry has changed but remains a nonfalsity theorem. MP, GEN, the strong compiler induction, exact candidate packing, and cofinal refutation padding are unchanged. This proves consistency equivalence for the new source.

\subsection{Physical distinction and initialization}
Reconstructing the decoder from the changed source and table is essential. An earlier inherited-allocation attempt kept stale roles and failed finite tests; that failure was not a proof that the two-row construction needed an extra state.

The new table is not an initial-state-preserving labelled-graph relabeling of either the old primary table or the public intake-199 allocation. The latter differs at 62 literal state/symbol keys; a forced rooted mapping also reaches a written-symbol/direction conflict, so this distinction does not rely on native label choices alone.

Its blank prefix is unchanged: neither modified row is executed before the certified root. Independent literal replay reaches the genuine source root after 95,714 transitions, with head $-8$, gap 1, 40 all-one spares, and $(N,T,W,A,D,H,X,Y,L)=(4,0,0,0,0,0,0,0,0)$. The complete sparse tape and dense relabeling are checked against the prefix receipt. The separate fresh arbitrary-natural runtime verification passes all 6,176 guarded cases, using 1,699,838 symbolic rules with minimum affine duration 6 and coverage of all 430 rows. Its exact table bindings and certificates are supplied with the secondary artifacts. The two machines in this report have independent physical certificates; their common state count is not used to transfer one proof to the other.

\section{Verification layers and reproducibility}
The archive separates the following claims. Passing one layer is not treated as proving the next.
\begin{enumerate}
\item \textbf{Literal machine.} Parse both transitions of every working state, verify all targets, count 215 working states and 430 rows, and check the native/normalized/canonical bijections. The canonical table in the appendix is generated from the machine file.
\item \textbf{Source identity.} Check the entire primitive word and complete guarded ordinal control graph against the calibrated primary source. This proves the all-natural source bisimulation, not binary execution.
\item \textbf{Binary runtime.} Recheck 6,176 guarded symbolic cases against the exact table and source. These use inductive scan/list lemmas and positive affine durations, rather than a fixed cutoff on register values.
\item \textbf{Blank prefix.} Independently replay all 716,307 transitions from zero tape and compare the entire final sparse tape. Decode its full frame and zero program counter; do not accept the first superficially well-formed interior frame.
\item \textbf{Source consistency.} Use the explicit numerical-calculus calibration, strong proof compiler, finite draining, and the new arbitrary-counter root theorem. This is what connects physical halting with ZFC inconsistency.
\item \textbf{Distinctness.} Check the graph-refinement certificate and exact source/table differences against the three pinned baselines. This excludes a mere relabeling relative to those inputs, not every possible previously published construction.
\end{enumerate}

The ordinary proof engine and its certificates are executable mathematical evidence. They are not a newly checked PA proof object, and this work does not supply a new exact-table Lean/Rocq endpoint. The inherited set-theoretic calibration is explicitly part of the provenance chain; this work does not claim to formalize it afresh. Runtime replay by the same checker is a fresh execution, not an independent second proof algorithm. The additional sparse-tape initializer and graph/source audits are independently implemented checks of their respective finite claims.

\subsection{What did not work toward 214}
A separate 214-state initialized core implements the original register source in all 6,176 unbounded primitive cases. It is \emph{not} a blank-start construction: starting it from any of its 214 existing states on blank tape leads to a verified outward ray. Consequently it cannot be substituted for the present 215-state machine. Context-sensitive startup repair remains open.

Three further 214-state tables passed an initially too-small macro sample. Broader admitted inputs exposed incorrect arithmetic, including an actual blank-run mismatch for one table. Those tables are not certified alternatives and are not included among the claimed results. This experience motivates explicitly distinguishing finite search filters from unbounded contracts and from a changed source that would require its own consistency proof.

\subsection{Provenance}
The public project is
\url{https://github.com/VladimirReshetnikov/ProveIt}, in
\file{Computability/BusyBeaver/Research/ZFC/}. The numerical dominated-row recovery traces back to commit
\file{7071e5041ee316d087661b658362feed159901ef}; the local-input permutation parent is commit
\file{d59b4418821c598c4703f0a60b1f7cbae4f3bcf6}. The previous composed K215 packet supplies the immediate source/runtime reference and reproduced calibration sections. Its two state tables, and the later helper-role variant, are pinned in the comparison evidence. Attribution and inherited licenses accompany the code rather than being replaced by a claim that this is an independent new calculus.

\clearpage\appendix
\section{Complete canonical transition table}
Start state 0, head coordinate 0, and an everywhere-zero tape are part of the definition. Working states are $0,\ldots,214$; $H$ denotes halt label 215, which has no outgoing transitions. An entry $wDq$ means write $w$, move one cell in direction $D$, and enter state $q$. The appendix is generated from the supplied canonical table.
\small\renewcommand{\arraystretch}{0.94}
\begin{longtable}{rll@{\qquad}rll}
\toprule State & Read 0 & Read 1 & State & Read 0 & Read 1\\\midrule
\endfirsthead
\toprule State & Read 0 & Read 1 & State & Read 0 & Read 1\\\midrule
\endhead
0 & $1R93$ & $1R0$ & 108 & $0R123$ & $1R8$\\
1 & $0L1$ & $1L2$ & 109 & $0R102$ & $1R124$\\
2 & $0L3$ & $1L2$ & 110 & $0R152$ & $1R190$\\
3 & $0L5$ & $1L2$ & 111 & $0L214$ & $0L96$\\
4 & $1R6$ & $0L4$ & 112 & $0R126$ & $1R105$\\
5 & $0L9$ & $1L2$ & 113 & $0R173$ & $1R153$\\
6 & $0R6$ & $0R10$ & 114 & $0R166$ & $1R184$\\
7 & $0R85$ & $1R194$ & 115 & $0R192$ & $1R166$\\
8 & $0R14$ & $1R13$ & 116 & $0R150$ & $1R148$\\
9 & $0L15$ & $1L2$ & 117 & $0R161$ & $1R205$\\
10 & $0L16$ & $0R10$ & 118 & $0L133$ & $0L1$\\
11 & $0L26$ & $1R131$ & 119 & $0R39$ & $1R73$\\
12 & $0R19$ & $1R20$ & 120 & $0R62$ & $1R134$\\
13 & $0L90$ & $1L31$ & 121 & $0R62$ & $1R135$\\
14 & $0L11$ & $0R22$ & 122 & $0R71$ & $0L16$\\
15 & $0L23$ & $1L2$ & 123 & $0R14$ & $1R134$\\
16 & $1L1$ & $0L1$ & 124 & $0R21$ & $1R136$\\
17 & $0R24$ & $0R32$ & 125 & $0R149$ & $1R181$\\
18 & $0R206$ & $1R205$ & 126 & $0R139$ & $1L20$\\
19 & $0R28$ & $1R29$ & 127 & $0R132$ & $1R183$\\
20 & $0R30$ & $1R31$ & 128 & $0R142$ & $1R200$\\
21 & $0R24$ & $0R60$ & 129 & $0R76$ & $1R179$\\
22 & $0R33$ & $1R22$ & 130 & $0R17$ & $1R180$\\
23 & $0L34$ & $1L2$ & 131 & $0R76$ & $1R138$\\
24 & $0L35$ & $0R0$ & 132 & $0R12$ & $1R143$\\
25 & $0R36$ & $1R25$ & 133 & $0R30$ & $1R70$\\
26 & $0R37$ & $1R38$ & 134 & $0R39$ & $1L126$\\
27 & $0R39$ & $0L40$ & 135 & $0R155$ & $0L16$\\
28 & $0R41$ & $1R42$ & 136 & $0R156$ & $0L16$\\
29 & $0R43$ & $1R44$ & 137 & $0R115$ & $1R88$\\
30 & $0R45$ & $1R46$ & 138 & $0R162$ & $1R141$\\
31 & $0R21$ & $1R47$ & 139 & $0R24$ & $0R48$\\
32 & $0R48$ & $1R32$ & 140 & $0R142$ & $1R183$\\
33 & $0R32$ & $1R33$ & 141 & $0R198$ & $1R95$\\
34 & $0L50$ & $1L2$ & 142 & $0R97$ & $1R160$\\
35 & $1R35$ & $1L1$ & 143 & $0R178$ & $1R167$\\
36 & $0R51$ & $1R36$ & 144 & $0R12$ & $1R160$\\
37 & $0R17$ & $1R211$ & 145 & $0R89$ & $1L167$\\
38 & $0R13$ & $1R53$ & 146 & $0R89$ & $1R159$\\
39 & $0L11$ & $0R36$ & 147 & $0L154$ & $1L164$\\
40 & $0L4$ & $1L54$ & 148 & $0R85$ & $1R89$\\
41 & $0R24$ & $0R25$ & 149 & $0R160$ & $1R12$\\
42 & $0R13$ & $0L16$ & 150 & $0R84$ & $1R145$\\
43 & $0R56$ & $1R57$ & 151 & $0R165$ & $1R169$\\
44 & $0R58$ & $1R59$ & 152 & $0L67$ & $1L143$\\
45 & $0R24$ & $0R55$ & 153 & $0R150$ & $1R207$\\
46 & $0R61$ & $1R62$ & 154 & $0R160$ & $1R116$\\
47 & $0R63$ & $0L16$ & 155 & $0R98$ & $1R98$\\
48 & $0R25$ & $1R48$ & 156 & $0R49$ & $0R32$\\
49 & $0L35$ & $0R64$ & 157 & $0R201$ & $1R128$\\
50 & $0L75$ & $1L2$ & 158 & $0L185$ & $1R97$\\
51 & $0R16$ & $1R51$ & 159 & $0R169$ & $1R182$\\
52 & $0R149$ & $1R145$ & 160 & $0R19$ & $1R172$\\
53 & $0R68$ & $1R69$ & 161 & $0R176$ & $1R166$\\
54 & $0L40$ & $1L54$ & 162 & $0R160$ & $1R97$\\
55 & $0R22$ & $1R55$ & 163 & $0R127$ & $1R165$\\
56 & $0L11$ & $0R51$ & 164 & $0R98$ & $1R163$\\
57 & $0R61$ & $0L16$ & 165 & $0L7$ & $1R144$\\
58 & $0R70$ & $1R71$ & 166 & $0L110$ & $1R176$\\
59 & $0R62$ & $1R72$ & 167 & $0R148$ & $1R175$\\
60 & $0R55$ & $1R60$ & 168 & $0R89$ & $1R157$\\
61 & $0R49$ & $0R25$ & 169 & $0R162$ & $1R145$\\
62 & $0L28$ & $1L73$ & 170 & $0R146$ & $1R137$\\
63 & $0L119$ & $1L30$ & 171 & $0R181$ & $1R181$\\
64 & $0R74$ & $1R64$ & 172 & $0R184$ & $1R124$\\
65 & $0R98$ & $1R152$ & 173 & $0R147$ & $1R166$\\
66 & $1R188$ & $1L54$ & 174 & $0R145$ & $1R170$\\
67 & $0R146$ & $1R147$ & 175 & $0R76$ & $1R61$\\
68 & $0R77$ & $1RH$ & 176 & $0R147$ & $1R89$\\
69 & $0L201$ & $1R31$ & 177 & $0R105$ & $1R97$\\
70 & $0R41$ & $1R63$ & 178 & $0R142$ & $1R85$\\
71 & $0L90$ & $1R79$ & 179 & $0R178$ & $1R76$\\
72 & $0R56$ & $1L80$ & 180 & $0R191$ & $0L214$\\
73 & $0R81$ & $1R82$ & 181 & $0R144$ & $1L151$\\
74 & $1R0$ & $1R64$ & 182 & $0R98$ & $1R52$\\
75 & $0R75$ & $1R83$ & 183 & $0R169$ & $1R52$\\
76 & $0L147$ & $1L143$ & 184 & $0R17$ & $1L133$\\
77 & $0R86$ & $1R87$ & 185 & $0R174$ & $1R186$\\
78 & $0R130$ & $1R12$ & 186 & $0R77$ & $1R118$\\
79 & $0R90$ & $1L134$ & 187 & $0R107$ & $1R126$\\
80 & $0L109$ & $1R57$ & 188 & $1L40$ & $1L111$\\
81 & $0R49$ & $0R33$ & 189 & $0R173$ & $1R160$\\
82 & $0R92$ & $0L57$ & 190 & $0R152$ & $1R117$\\
83 & $0R94$ & $1R202$ & 191 & $0R11$ & $1R99$\\
84 & $0R97$ & $1R12$ & 192 & $0R196$ & $1R100$\\
85 & $0R89$ & $1R98$ & 193 & $0R148$ & $1R203$\\
86 & $0R213$ & $1R21$ & 194 & $0R158$ & $1R195$\\
87 & $0R118$ & $1R101$ & 195 & $0R178$ & $1R189$\\
88 & $0R152$ & $1R114$ & 196 & $0R65$ & $1R7$\\
89 & $0R103$ & $1R104$ & 197 & $0R152$ & $1R129$\\
90 & $0R81$ & $1R80$ & 198 & $0R69$ & $1R212$\\
91 & $0R102$ & $1R177$ & 199 & $0R154$ & $1R125$\\
92 & $0R106$ & $1R107$ & 200 & $0R165$ & $1R185$\\
93 & $0R66$ & $0R0$ & 201 & $0R91$ & $1R89$\\
94 & $0R27$ & $1R108$ & 202 & $0R109$ & $1R26$\\
95 & $0R186$ & $0L214$ & 203 & $0R210$ & $1R110$\\
96 & $1L111$ & $1L96$ & 204 & $0R152$ & $1R171$\\
97 & $0R19$ & $1R112$ & 205 & $0R204$ & $1R161$\\
98 & $0R61$ & $1R61$ & 206 & $0R142$ & $1R140$\\
99 & $0R201$ & $1R208$ & 207 & $0R199$ & $1R18$\\
100 & $0R67$ & $1R193$ & 208 & $0R28$ & $1R158$\\
101 & $0R106$ & $0L214$ & 209 & $0R164$ & $1R197$\\
102 & $0R119$ & $1R120$ & 210 & $0R209$ & $1R113$\\
103 & $0R72$ & $1R59$ & 211 & $0R78$ & $0L214$\\
104 & $0R58$ & $1R121$ & 212 & $0R168$ & $1R187$\\
105 & $0R21$ & $1R122$ & 213 & $0R45$ & $1R101$\\
106 & $0R21$ & $0L10$ & 214 & $0L4$ & $0L214$\\
107 & $0R62$ & $1R42$ & & &\\
\bottomrule
\end{longtable}
\renewcommand{\arraystretch}{1}\normalsize

\end{document}
